% This is file JFM2esam.tex
% first release v1.0, 20th October 1996
%       release v1.01, 29th October 1996
%       release v1.1, 25th June 1997
%       release v2.0, 27th July 2004
%       release v3.0, 16th July 2014
%   (based on JFMsampl.tex v1.3 for LaTeX2.09)
% Copyright (C) 1996, 1997, 2014 Cambridge University Press

\documentclass{jpp}
\usepackage{graphicx}
\usepackage{epstopdf, epsfig}
\usepackage{hyperref}
\usepackage{amsmath,amsfonts,amssymb}
\usepackage{subfigure}

\newtheorem{theorem}{Theorem}
\newtheorem{lemma}{Lemma}
\newtheorem{corollary}{Corollary}
\newtheorem{remark}{Remark}


%\shorttitle{Short title Short title Short title}
\shortauthor{A. Author One, B. Author Two and C. Author Three}

\title{ Capture threshold for electron interacting with circularly polarized waves in parallel static electric fields}

\author{A. Author One\aff{1}\corresp{\email{abcd.xyz@mnop.edu}}, B. Author Two\aff{2} and C. Author Three\aff{3}}

\affiliation{\aff{1}University of Science and Technology of China
\aff{2}Affiliation Two Department Name City Name Country Name
\aff{3}Affiliation Three Department Name City Name Country Name}

\begin{document}

\maketitle

\begin{abstract}
%The electron's parallel  acceleration under parallel electric field would be convert to cyclotron velocity increasing while keep in resonant condition when  captured by the circularly waves.     The capture threshold of electron interacting with circularly waves in parallel static electric fields has been studied through the nonlinear Schrodinger-type equation,  numerical results and analyzed expression are compared directly , which show strongly agreement. Capture condition with R/L- waves are also discussed
This study investigates the threshold for resonant capture of electrons interacting with circularly polarized electromagnetic waves in a static electric field parallel to the background magnetic field. Captured electrons can channel energy supplied by the static electric field into perpendicular cyclotron motion while maintaining resonance. An analytical expression for the critical wave amplitude required for capture is derived using a nonlinear Schrödinger-type equation, characterizing the transition between passage through resonance and sustained resonant capture. The analytical prediction is compared directly with numerical results, showing strongly agreement. Capture conditions for right- and left-hand circularly polarized waves are also analyzed to examine the influence of wave polarization on resonant electron dynamics.
\end{abstract}

\keywords{Autoresonance,Capture Threshold, Anomalous Doppler resonance, Normal Doppler resonance}

%\section*{Nomenclature}
%
%\noindent \begin{tabular}{@{}l@{\qquad}l@{}}
%$a$ &  definition text\\
%$a,b$ &  definition text definition text\\
%APU &  definition text definition text \\
%$ci$ &  definition text definition text    \\
%$CT$ &  definition text \\
%CAS &  definition text definition text \\
%$E$ &  definition text \\
%$D$ &  definition text  \\
%$f$ &  definition text   \\
%$Fg$ &  definition text\\
%$F$ &  definition text  \\
%FCOM &  definition text     \\
%\end{tabular}

%\section*{Greek Symbol}
%
%\noindent\begin{tabular}{@{}l@{\qquad}p{28pc}@{}}
%$\alpha$ &  definition text  \\
%$\beta$ &  definition text definition text  \\
%$\gamma$ &  definition text definition text  \\
%$\delta$ &  definition text definition text  \\
%$\xi$ &  definition text definition text \\
%$\phi$ &  definition text definition text definition text definition text definition text definition
%text definition text definition text definition text \\
%$\psi$ &  definition text definition text \\
%\end{tabular}

\section{Introduction}



\subsection{Second-Order Heading}

This is an example of dummy text. This is an example of dummy text. This is an example
of dummy text. This is an example of dummy text. This is an example of dummy text. This
is an example of dummy text. This is an example of dummy text. This is an example of
dummy text. This is an example of dummy text. This is an example of dummy text. This is
an example of dummy text.\footnote{This is an example of first footnote.}

\subsection{Lists}

Here is an example of a numbered list:

\begin{enumerate}
\item[1.] List entry number 1 List entry number 1 List entry number 1 List entry number 1 List
entry number 1 List entry number 1,
\item[2.] List entry number 2, List entry number 2, List entry number 2, List entry number 2, List
entry number 2, List entry number 2,
\item[3.] List entry number 3,
\item[4.] List entry number 4, and
\item[5.] List entry number 5.
\end{enumerate}

The following is an example of an \textit{itemized/bulleted} list.

\begin{itemize}
\item This is an example of bulleted listing.
\item This is an example of bulleted listing. This is an example of bulleted listing.
\item This is an example of bulleted listing. This is an example of bulleted listing. This is an
example of bulleted listing.
\item This is an example of bulleted listing. This is an example of bulleted listing. This is an
example of bulleted listing. This is an example of bulleted listing.
\item This is an example of bulleted listing.
\end{itemize}

This is an example of dummy text. This is an example of dummy text. This is an example
of dummy text. This is an example of dummy text. This is an example of dummy text. This is
an example of dummy text. This is an example of dummy text. This is an example of dummy
text. This is an example of dummy text. This is an example of dummy text. This is an example
of dummy text. This is an example of dummy text. This is an example of dummy text. This is
an example of dummy text. This is an example of dummy text. This is an example of dummy
text. This is an example of dummy text. This is an example of dummy text. This is an example
of dummy text. This is an example of dummy text. This is an example of dummy text. This is
an example of dummy text. This is an example of dummy text. This is an example of dummy
text.

\begin{itemize}
\item This is an example of bulleted listing.
\item This is an example of bulleted listing. This is an example of bulleted listing.
\item This is an example of bulleted listing. This is an example of bulleted listing. This is an
example of bulleted listing. It contains another list nested inside it. The inner list is an
\textit{enumerated} list.

\begin{enumerate}
\item[1.] This is the first item of an enumerated list that is nested within the itemized list.
\item[2.] This is the second item of the inner list. LATEX allows you to nest lists deeper than
you really should.
\end{enumerate}
\item[]This is the rest of the second item of the outer list. It is no more interesting than any other
part of the item.
\item This is the third item of the list.
\end{itemize}

The default formatting of numbering can be changed as (1), (2), (3). . . , for example,

\begin{enumerate}[(5)]
\item[(1)] List entry one.
\item[(2)] List entry two.
\item[(3)] List entry three.
\item[(4)] List entry four.
\item[(5)] List entry five.
\end{enumerate}

The default formatting of numbering can be changed as (i), (ii), (iii). . . , for example,

\begin{enumerate}[(iii)]
\item[(i)] List entry one.
\item[(ii)] List entry two.
\item[(iii)] List entry three.
\item[(iv)] List entry four.
\item[(v)] List entry five.
\end{enumerate}

The default formatting of numbering can be changed as (a), (b), (c). . . , for example,

\begin{enumerate}[(e)]
\item[(a)] List entry.
\item[(b)] List entry two.
\item[(c)] List entry three.
\item[(d)] List entry four.
\item[(e)] List entry five.
\end{enumerate}

%\begin{bottomfigure}
%\vskip1pc
%\centering{\includegraphics{dummy-image}}
%\caption{This is an example of a figure caption. (a) This is a description for part a,
%and (b) this is a description for part b.}
%\end{bottomfigure}

The default formatting of numbering can be changed as I, II, III. . . , for example,

\begin{enumerate}[III.]
\item[I.] List entry.
\item[II.] List entry two.
\item[III.] List entry three.
\item[IV.] List entry four.
\item[V.] List entry five.
\end{enumerate}

\subsection{Here is an example of extract}

\begin{quotation}
\noindent This is an example text of quote or extract. This is an example text of quote or extract.
This is an example text of quote or extract. This is an example text of quote or extract.
This is an example text of quote or extract. This is an example text of quote or extract.
This is an example text of quote or extract. This is an example text of quote or extract.
\end{quotation}

This is dummy text. This is dummy text. This is dummy text. This is dummy text. This is
dummy text.

\subsection{Equations}

Example of a single-line equation 1.1.
\begin{equation}
a = b\ \text{((Single Equation Numbered)) } \hfill. . .
\end{equation}


Equations can also be multiple lines as shown in Equations 1.2 and 1.3.
\begin{align}
c &= 0\ \text{((Multiple Lines, Numbered))}\hfill . . .\\
ac &= 0\ \text{((Multiple Lines, Numbered))} \hfill . . .
\end{align}

Unnumbered display equations single line:
\begin{equation*}
\Gamma \psi' = x'' + y^2 +z _i^n
\end{equation*}

Unnumbered display equations multiple lines:
\begin{align*}
c &= 0\ \text{((Multiple Lines, Unnumbered))}\hfill . . .\\
ac &= 0\ \text{((Multiple Lines, Unnumbered))} \hfill . . .
\end{align*}

\subsubsection{Third-Order Heading}

This is dummy text. This is dummy text. This is dummy text. This is dummy text. This is
dummy text.

This is dummy text. This is dummy text. This is dummy text. This is dummy text. This is
dummy text. This is dummy text. This is dummy text. This is dummy text. This is dummy
text. This is dummy text. This is dummy text. This is dummy text. This is dummy text. This
is dummy text. This is dummy text. This is dummy text. This is dummy text. This is dummy
text. This is dummy text. This is dummy text in Fig. 1.

%\begin{figure}
%\centering{\includegraphics{dummy-image}}
%\caption{This is an example of a figure caption. (a) This is a description for part a,
%and (b) this is a description for part b.}
%\end{figure}

This is dummy text. This is dummy text. This is dummy text. This is dummy text. This is
dummy text. This is dummy text. This is dummy text. This is dummy text. This is dummy
text. This is dummy text. This is dummy text. This is dummy text. This is dummy text. This
is dummy text. This is dummy text. This is dummy text. This is dummy text. This is dummy
text. This is dummy text. This is dummy text. This is dummy text. This is dummy text. This
is dummy text. This is dummy text. This is dummy text.


\begin{table}
\begin{center}
\caption{Table caption}
{\begin{tabular}{lcccc}
\textbf{Year} & \textbf{Single outlet} &  \textbf{Small multiple}$^{\bf a}$ & \textbf{Large multiple} & \textbf{Total}\\[6pt]
1982$^{b}$ & 98 & 129 & 620 & 847\\
1987 & 138 & 176 & 1000 & 1314\\
1991 & 173 & 248 & 1230 & 1651\\
1998 & 200 & 300 & 1500 & 2000\\[6pt]
\end{tabular}}{\par This is an example of unnumbered tablenote\\
$^{a}$This is an example of first numbered tablenote\\
$^{b}$This is an example of second numbered tablenote\\}
\end{center}
\end{table}

This is dummy text. This is dummy text. This is dummy text. This is dummy text. This is
dummy text. This is dummy text. This is dummy text. This is dummy text. This is dummy
text. This is dummy text. This is dummy text in Fig. 2.

This is dummy text. This is dummy text. This is dummy text. This is dummy text. This is
dummy text. This is dummy text. This is dummy text. This is dummy text. This is dummy
text. This is dummy text. This is dummy text. This is dummy text. This is dummy text. This
is dummy text. This is dummy text. This is dummy text. This is dummy text. This is dummy
text. This is dummy text. This is dummy text. This is dummy text. This is dummy text. This
is dummy text. This is dummy text. This is dummy text in Table 1.

Figure 3 shows the dummy text. This is dummy text. This is dummy text. This is dummy
text. This is dummy text. This is dummy text. This is dummy text. This is dummy text. This
is dummy text. This is dummy text. This is dummy text. This is dummy text. This is dummy
text. This is dummy text. This is dummy text. This is dummy text. This is dummy text. This
is dummy text. This is dummy text. This is dummy text. This is dummy text. This is dummy
text. This is dummy text. This is dummy text in Table 2.

This is dummy text.\footnote{This is an example of second footnote.} This is dummy text. This is dummy text. This is dummy text. This
is dummy text. This is dummy text. This is dummy text. This is dummy text. This is dummy
text. This is dummy text. This is dummy text. This is dummy text. This is dummy text. This
is dummy text. This is dummy text. This is dummy text. This is dummy text. This is dummy
text. This is dummy text. This is dummy text. This is dummy text. This is dummy text. This
is dummy text in Table 2.

%\begin{figure}
%\centering{\includegraphics{dummy-image}}
%\caption{This is an example of a figure caption. This is an example of a figure caption. This is an example of a
%figure caption. This is an example of a figure caption. This is an example of a figure caption. This is an
%example of a figure caption. (a) This is a description for part a, and (b) this is a description for part b.}
%\end{figure}

\begin{table}
\begin{center}
\caption{Table caption table caption table caption table
caption table caption table caption~table}
{\begin{tabular}{lcccc}
\textbf{Column 1} & \textbf{Column 2}  & \textbf{Column 3}  & \textbf{Column 4}  & \textbf{Column 5}\\[6pt]
460  & 2.12  & 1.04  & 8.63  & 3.48\\
476  & 2.24  & 0.94  & 9.22  & 3.73\\
676  & 3.71  & 1.67  & 16.36  & 6.00\\
572  & 2.86  & 1.29  & 11.90  & 4.56\\[6pt]
\end{tabular}}{\par $^{a}$First Table Footnote;\\
$^{b}$TSecond Table Footnote.\\}
\end{center}
\end{table}

This is an example of theorem.

\begin{theorem}
This is an example of enunciation. This is an example of enunciation. This is an
example of enunciation. This is an example of enunciation. This is an example of enunciation.
This is an example of enunciation.
\begin{equation*}
AB = \sum_{i} a_{i} b_{i}
\end{equation*}
\end{theorem}

\section{FIRST-ORDER HEADING}
\paragraph{Fourth-Order Heading}

This is an example of Lemma
\begin{lemma}
This is an example of enunciation. This is an example of enunciation. This is an
example of enunciation. This is an example of enunciation. This is an example of enunciation.
This is an example of enunciation.
\begin{equation*}
AB = \sum_{i} a_{i} b_{i}
\end{equation*}
\end{lemma}

\begin{theorem}
This is an example of enunciation. This is an example of enunciation. This is an
example of enunciation. This is an example of enunciation. This is an example of enunciation.
This is an example of enunciation.
\begin{equation*}
AB = \sum_{i} a_{i} b_{i}
\end{equation*}
\end{theorem}

\begin{theorem}[Math Detail]
This is an example of enunciation. This is an example of enunciation. This is an
example of enunciation. This is an example of enunciation. This is an example of enunciation.
This is an example of enunciation.
\begin{equation*}
AB = \sum_{i} a_{i} b_{i}
\end{equation*}
\end{theorem}

This is an example of Remarks.

\begin{remark}
This is an example of enunciation. This is an example of enunciation. This is an
example of enunciation. This is an example of enunciation. This is an example of enunciation.
This is an example of enunciation.
\begin{equation*}
AB = \sum_{i} a_{i} b_{i}
\end{equation*}
\end{remark}

This is dummy text. This is dummy text. This is dummy text. This is dummy text. This is
dummy text. This is dummy text. This is dummy text. This is dummy text. This is dummy
text. This is dummy text. This is dummy text. This is dummy text. This is dummy text. This
is dummy text. This is dummy text. This is dummy text. This is dummy text. This is dummy
text. This is dummy text.

\section*{Acknowledgements}

This is an example of acknowledgement$^{(1,2)}$. This is an example of acknowledgement. This
is an example of acknowledgement. This is an example of acknowledgement$^{(3)}$. This is an
example of acknowledgement. This is an example of acknowledgement. This is an example
of acknowledgement. This is an example of acknowledgement$^{(4)}$.\cite{RN2131}


\makeatletter
\def\fps@table{h}
\def\fps@figure{h}
\makeatother
\bibliographystyle{jpp}
\bibliography{ustc4}
\section*{Appendix}

\section*{First order heading}

\begin{table}
\begin{center}
\caption{Table caption table caption table caption table caption table caption table
caption table caption table caption table caption table caption table
caption table caption table caption table caption}
{\tabcolsep1pc\begin{tabular}{@{}lccccc@{}}
& & \multicolumn{2}{c}{\textbf{Span Col 3 \& 4}} & \multicolumn{2}{c}{\textbf{Span Col 5 \& 6}}\\
\textbf{Column} & \textbf{Column} & \textbf{Column} & \textbf{Column} & \textbf{Column} & \textbf{Column}\\
\textbf{One} & \textbf{Two} & \textbf{Three} & \textbf{Four} & \textbf{Five} & \textbf{Six}\\[6pt]
6-31+G(d,p) & 460 & 2.12 & 1.04 & 8.63 & 3.48\\
6-31++G(d,p) & 476 & 2.24 & 0.94 & 9.22 & 3.73\\
6-31++G(df,p) & 676 & 3.71 & 1.67 & 16.36 & 6.00\\
6-311+G(d,p) & 572 & 2.86 & 1.29 & 11.90 & 4.56\\
6-311++G(3df,2p) & 1,060 & 7.54 & 3.25 & 36.70 & 11.45\\[6pt]
\end{tabular}}{\par$^{a)}$Table Footnote\\
$^{b)}$Second Table Footnote.}
\end{center}
\end{table}
%
%\begin{figure}
%\begin{minipage}{.4\textwidth}
%\centering
%\includegraphics[scale=0.4]{dummy-image}
%\caption{This is an example of a figure caption. (a) This is a description for part a,
%and (b) this is a description for part b.}
%\end{minipage}\hskip3pc
%\begin{minipage}{.5\textwidth}
%\centering
%\includegraphics[scale=0.6]{dummy-image}
%\caption{This is an example of a figure caption. (a) This is a description for part a,
%and (b) this is a description for part b.}
%\end{minipage}
%\end{figure}



\end{document}
